% ============================================================================
% DÉFINITION DU COURS 07 : ÉLECTRONIQUE — ACQUISITION ET TRAITEMENT DU SIGNAL
% ============================================================================

\titlehead{Thomas Lusseau - SII en ATS}
\title[Électronique]{Acquisition et traitement du signal}
\author[TL]{Thomas Lusseau}
\date{\today}

\TLCourseCoverSetup
  {\repStyle/FondPoly.png}
  {\repStyle/FondPoly.png}
  {Cycle 4 : Acquérir, conditionner et traiter l'information}
  {Acquisition et traitement\\ du signal}
  {Lycée Robert Doisneau - ATS}
  {Thomas Lusseau}

\TLCourseCoverContents{%
  \TLCourseContentsLine{7.1}{Chaîne d'acquisition}{sec:elec-chaine}
  \TLCourseContentsLine{7.2}{Signaux analogiques et numériques}{sec:elec-signaux}
  \TLCourseContentsLine{7.3}{Capteurs résistifs et ponts de mesure}{sec:elec-capteurs}
  \TLCourseContentsLine{7.4}{Amplificateur opérationnel}{sec:elec-aop}
  \TLCourseContentsLine{7.5}{Filtrage analogique}{sec:elec-filtrage}
  \TLCourseContentsLine{7.6}{Échantillonnage}{sec:elec-echantillonnage}
  \TLCourseContentsLine{7.7}{Conversion analogique-numérique}{sec:elec-can}
  \TLCourseContentsLine{7.8}{Conditionnement et qualité de mesure}{sec:elec-conditionnement}
}

\TLCourseFooter
\TLCourseCover
\markboth{Acquisition et traitement du signal}{Acquisition et traitement du signal}
\TLSetReflectionImage{\repStyle/FondPoly.png}

\livrettrue
\collefalse
\setcounter{chapter}{6}
\chapter*{Acquisition et traitement du signal}
\addcontentsline{toc}{chapter}{Acquisition et traitement du signal}

% ============================================================================
% COURS 07 — ACQUISITION ET TRAITEMENT DE L'INFORMATION
% Reprise fidèle du document 11-AcquisitionTraitementSignal2022.docx
% ============================================================================

\section*{Objectifs du cours}

Ce cours donne les démarches et les outils nécessaires à l'étude des chaînes d'information.

\subsection*{Savoirs}

À l'issue du cours, il faut connaître :
\begin{itemize}
  \item la structure d'une chaîne d'acquisition : capteur, conditionneur, amplificateur, filtre, convertisseur analogique-numérique, unité de traitement et convertisseur numérique-analogique ;
  \item la différence entre résolution et précision ;
  \item les impédances et admittances des dipôles de base ;
  \item les différents types de filtres ;
  \item la méthode permettant de déterminer rapidement le comportement d'un filtre aux basses et aux hautes fréquences ;
  \item le théorème de Shannon-Nyquist ;
  \item le filtre anti-repliement et son utilité ;
  \item la définition du quantum d'un CAN et d'un CNA ;
  \item les méthodes permettant de déterminer la résolution d'un CAN ou d'un CNA ;
  \item les bases 2, 10 et 16 et les méthodes de changement de base ;
  \item le codage en complément à deux et ses particularités ;
  \item les opérations simples en base 2 ;
  \item la notion de temps de conversion d'un CAN et sa relation avec la période d'échantillonnage.
\end{itemize}

\subsection*{Savoir-faire}

Il faut savoir :
\begin{itemize}
  \item identifier les capteurs utilisés dans la chaîne d'information d'un système ;
  \item identifier les constituants matériels de la chaîne d'information et les fonctions techniques réalisées ;
  \item définir la nature des informations d'entrée et de sortie d'un capteur ;
  \item justifier le choix d'un capteur au regard du cahier des charges ;
  \item extraire d'une documentation les valeurs numériques caractéristiques des solutions retenues ;
  \item analyser et mettre en œuvre une association de fonctions électroniques réalisant le conditionnement d'un signal ;
  \item décrire et représenter les évolutions temporelles et fréquentielles des signaux ;
  \item tracer le spectre d'un signal à partir de son expression temporelle et réaliser l'opération inverse ;
  \item déterminer la fonction de transfert d'un filtre passif ;
  \item prévoir le comportement d'un filtre aux basses et hautes fréquences ;
  \item tracer le spectre de sortie d'un filtre à partir du spectre d'entrée ;
  \item tracer le gabarit d'un filtre et déterminer l'ordre minimal associé ;
  \item déterminer la fréquence minimale d'échantillonnage ;
  \item déterminer le quantum, la résolution et le nombre de bits d'un CAN ou d'un CNA ;
  \item déterminer le gain d'un amplificateur permettant d'obtenir une résolution donnée ;
  \item coder une valeur en bases 2, 10 et 16 ainsi qu'en complément à deux.
\end{itemize}

\section{Chaîne d'acquisition}\label{sec:elec-chaine}

\subsection{Introduction}

Tous les systèmes pluritechniques ont besoin de contrôler de nombreux paramètres physiques : température, force, vitesse, position, déplacement, pression ou luminosité. À chacune de ces grandeurs peuvent correspondre un ou plusieurs types de capteurs, fondés sur des phénomènes physiques différents : variation de résistance, variation d'induction magnétique, effet Hall, effet piézoélectrique, etc.

Le capteur fournit un signal électrique — tension, courant ou variation d'impédance — image de la grandeur physique mesurée. Des circuits électroniques mettent ensuite en forme ce signal avant son exploitation par une unité de traitement : microcontrôleur, automate ou ordinateur.

\TLChaineAcquisition

La chaîne d'acquisition peut comporter les fonctions suivantes :
\begin{center}
\textbf{capteur et conditionneur $\rightarrow$ amplification $\rightarrow$ filtrage $\rightarrow$ CAN $\rightarrow$ unité de traitement $\rightarrow$ CNA.}
\end{center}

L'information fournie à l'unité de traitement peut se présenter sous trois formes :
\begin{itemize}
  \item signal analogique ;
  \item signal logique ;
  \item signal numérique.
\end{itemize}

\subsection{Nature du signal : analogique, logique et numérique}

Toutes les grandeurs qui nous entourent sont initialement analogiques : leur valeur est définie à chaque instant et peut prendre une infinité de valeurs. Cette forme est incompatible avec les seuls états « 0 » et « 1 » compris par une unité numérique.

Deux conversions sont alors possibles :
\begin{itemize}
  \item convertir l'information en un signal logique, ou signal tout ou rien ;
  \item convertir l'information en un signal numérique.
\end{itemize}

\subsubsection{Signal logique ou tout ou rien}

Un signal logique est obtenu en comparant le signal analogique à un seuil. Selon que la grandeur est inférieure ou supérieure à ce seuil, la sortie vaut « 0 » ou « 1 ». Cette transformation rend l'information directement exploitable par l'unité de traitement, mais elle fait perdre les détails relatifs aux variations de la grandeur analogique.

\paragraph{Exemple du destructeur d'aiguilles.}
Le destructeur d'aiguilles utilise un capteur infrarouge constitué d'une diode émettrice et d'une photodiode réceptrice. La présence ou l'absence de l'embase est interprétée par le microcontrôleur sous la forme d'un signal tout ou rien.

\subsubsection{Signal numérique}

Pour conserver davantage d'information qu'avec un signal logique, chaque valeur du signal analogique est codée par une suite de zéros et de uns appelée mot binaire.

Le signal est prélevé à intervalles réguliers : c'est l'échantillonnage. Cette opération entraîne une perte d'information, qui peut être limitée en choisissant une fréquence d'échantillonnage suffisante. Un compromis reste nécessaire : une fréquence trop élevée augmente la quantité de données à traiter et à mémoriser.

\subsection{Fonction « Traiter »}

La fonction « Traiter » dépend de la nature des informations :
\begin{itemize}
  \item les systèmes logiques traitent des informations bit à bit à l'aide de tables de vérité ou de graphes d'états ;
  \item les systèmes numériques traitent des mots binaires par programmation à partir d'algorithmes ou d'algorigrammes.
\end{itemize}

\section{Généralités sur les capteurs}\label{sec:elec-capteurs}

\subsection{Structure d'un capteur et vocabulaire utilisé}

Un capteur convertit une grandeur physique, appelée \emph{mesurande}, en un signal électrique image de ses variations. On parle aussi de transducteur. Dans la pratique, plusieurs conversions successives peuvent être nécessaires.

Le vocabulaire utilisé est le suivant :
\begin{description}
  \item[Mesurande :] grandeur physique que l'on souhaite mesurer : température, pression, force, position, etc.
  \item[Corps d'épreuve :] élément réagissant à la grandeur physique et fournissant une mesurande secondaire exploitable.
  \item[Transducteur :] élément lié au corps d'épreuve et transformant la mesurande secondaire en tension, courant ou variation d'impédance.
  \item[Conditionneur :] circuit électronique délivrant un signal exploitable à partir du signal brut du capteur. Il peut comporter une amplification, une mise en forme, un filtrage, un CAN ou un CNA.
\end{description}

\subsection{Différents capteurs}

Les systèmes mettent en œuvre de nombreux capteurs :
\begin{itemize}
  \item déplacements : potentiomètres résistifs, capteurs inductifs ou capacitifs, capteurs de proximité ;
  \item efforts et déformations : jauges d'extensiométrie, capteurs piézoélectriques ;
  \item courant : capteurs à effet Hall ;
  \item vitesse et position : génératrices tachymétriques, codeurs incrémentaux ou absolus, potentiomètres, résolveurs, tachymètres optiques ;
  \item pression, humidité et couple ;
  \item température : thermocouples, thermistances CTN et CTP, sondes Pt100 ;
  \item optique : photorésistances et fibres optiques.
\end{itemize}

\subsubsection{Classification par la nature du signal de sortie}

On distingue les capteurs analogiques, logiques et numériques. Un capteur analogique délivre un signal variant continûment. Un capteur intelligent peut intégrer toute la chaîne d'acquisition et communiquer directement avec un réseau ou un bus de terrain.

\subsubsection{Classification par le principe de transduction}

On distingue :
\begin{itemize}
  \item les capteurs passifs, pour lesquels la mesurande agit sur la résistance, l'inductance ou la capacité du transducteur ;
  \item les capteurs actifs, qui produisent directement une tension, un courant ou une charge électrique.
\end{itemize}

Quelques effets physiques classiques sont résumés ci-dessous.

\begin{center}
\begin{tabular}{lll}
\toprule
Mesurande & Effet utilisé & Grandeur électrique obtenue \\
\midrule
Température & thermoélectricité & tension \\
Flux optique & photoémission & courant \\
Température & pyroélectricité & charge \\
Force, pression, accélération & piézoélectricité & charge \\
Position & effet Hall & tension \\
Vitesse & induction électromagnétique & tension \\
\bottomrule
\end{tabular}
\end{center}

\subsection{Caractéristiques d'un capteur}

Le choix d'un capteur résulte d'un compromis entre plusieurs caractéristiques :
\begin{itemize}
  \item étendue de mesure ;
  \item résolution ;
  \item sensibilité ;
  \item précision ;
  \item rapidité ;
  \item coût ;
  \item bande passante.
\end{itemize}

\begin{description}
  \item[Étendue de mesure :] valeurs extrêmes accessibles à la grandeur mesurée.
  \item[Résolution :] plus petite variation de la grandeur d'entrée que le capteur peut déceler.
  \item[Sensibilité :] variation du signal de sortie rapportée à une variation de la grandeur mesurée ; elle correspond localement à la dérivée de la caractéristique entrée-sortie.
  \item[Précision :] aptitude du capteur à fournir une valeur proche de la valeur vraie.
  \item[Rapidité :] temps nécessaire pour que la sortie réagisse à une variation de l'entrée.
  \item[Bande passante :] intervalle de fréquences dans lequel le capteur fournit une mesure exploitable.
\end{description}

\begin{center}
\fbox{\parbox{0.85\linewidth}{\centering\textbf{Il faut bien distinguer la résolution, qui caractérise la plus petite variation détectable, et la précision, qui caractérise la proximité de la valeur mesurée avec la valeur vraie.}}}
\end{center}

\subsection{Conditionnement}

Les signaux issus des capteurs sont souvent faibles, parfois de l'ordre de quelques millivolts. Ils doivent être amplifiés avant d'être transmis à l'unité de traitement. Des convertisseurs peuvent aussi modifier la nature du signal : courant-tension, fréquence-tension, charge-tension, etc.

\subsection{Capteurs et conditionneurs classiques}

Quelques associations usuelles doivent être connues :
\begin{itemize}
  \item jauges de contraintes associées à un pont de Wheatstone ;
  \item accéléromètres utilisés pour mesurer une accélération ou détecter le sens d'un mouvement ;
  \item sondes de température Pt100, CTN ou CTP, parfois associées à des résistances de linéarisation ;
  \item capteurs de courant à effet Hall, utiles pour réaliser une boucle d'asservissement en courant.
\end{itemize}

\subsubsection{Application : accéléromètre ADXL335}

On utilise un accéléromètre ADXL335 pour mesurer une accélération maximale $\gamma_{nc,\max}=\SI{1.2}{\metre\per\second\squared}$. Le composant est alimenté sous $\SI{3}{\volt}$ et consomme $\SI{350}{\micro\ampere}$. Sa sensibilité est de $\SI{300}{\milli\volt\per g}$, sa plage de linéarité est de $\pm 3g$ et sa sortie vaut $\SI{1.5}{\volt}$ pour une accélération nulle.

\begin{enumerate}
  \item Déterminer la plage de tension correspondant au domaine de linéarité.
  \item Convertir ce domaine de linéarité en $\si{\metre\per\second\squared}$.
  \item Déterminer la tension mesurée pour une accélération de $\SI{1.2}{\metre\per\second\squared}$.
  \item Calculer la puissance consommée par le dispositif.
\end{enumerate}

\subsubsection{Application : mesure de température}

On souhaite mesurer la température d'une pièce entre $\SI{0}{\degreeCelsius}$ et $\SI{50}{\degreeCelsius}$ à l'aide d'un capteur donné pour une plage de $-\SI{50}{\degreeCelsius}$ à $\SI{100}{\degreeCelsius}$.

\begin{enumerate}
  \item Définir les caractéristiques nécessaires au choix de ce capteur.
  \item Préciser les critères permettant de comparer plusieurs solutions.
  \item À partir de la loi $R(T)$ fournie par la documentation du capteur, déterminer s'il s'agit d'une thermistance CTN ou CTP.
\end{enumerate}

\section{Outils mathématiques pour la chaîne d'information : électrocinétique}\label{sec:elec-electrocinetique}

\subsection{Généralités}

L'énergie électrique nécessaire au fonctionnement d'un circuit est fournie par des sources d'électricité appelées générateurs. Elle est disponible sous deux formes :
\begin{itemize}
  \item les sources continues, ou DC : piles, accumulateurs, panneaux solaires ;
  \item les sources alternatives, ou AC : centrales électriques, éoliennes, alternateurs.
\end{itemize}

L'énergie électrique est principalement produite sous forme alternative à partir de centrales qui transforment une énergie nucléaire, thermique, hydraulique, éolienne ou marémotrice. Son unité est le joule. En électricité, on utilise également le kilowattheure :
\[
\SI{1}{\kilo\watt\hour}=\SI{3.6}{\mega\joule}.
\]

\subsection{Électrocinétique}

L'électrocinétique étudie la circulation des charges électriques dans un circuit ou réseau électrique. Un circuit élémentaire est constitué de conducteurs métalliques et de composants électriques. Il comporte au minimum un générateur, un récepteur et les conducteurs qui les relient.

L'analogie hydraulique permet d'interpréter le fonctionnement : une pompe crée une différence de pression et met le fluide en mouvement ; de même, un générateur crée une différence de potentiel et provoque la circulation des charges.

\subsection{Charge et courant}

L'intensité électrique est définie comme le débit de charges traversant une section du conducteur :
\[
i(t)=\frac{\mathrm{d}q(t)}{\mathrm{d}t}.
\]

Le courant s'exprime en ampères et la charge en coulombs. La charge élémentaire de l'électron vaut $e=-1{,}6\times10^{-19}\,\si{\coulomb}$. Le sens conventionnel du courant est opposé au déplacement des électrons dans un conducteur métallique.

Un courant se mesure avec un ampèremètre placé en série dans la branche étudiée. Pour un régime continu, l'intensité est constante. Pour un régime variable, elle dépend du temps.

\subsection{Potentiel et différence de potentiel}

Chaque point d'un circuit possède un potentiel électrique $V$, défini par rapport à un point de référence choisi arbitrairement, appelé masse ou potentiel nul.

La tension entre deux points $A$ et $B$ est la différence de potentiel :
\[
U_{AB}=V_A-V_B.
\]

Une tension se mesure avec un voltmètre placé en dérivation entre les deux points. L'orientation de la flèche fixe l'ordre des potentiels dans la relation.

\subsection{Dipôles}

Un dipôle est un élément électrique comportant deux bornes. Son comportement est caractérisé par une relation entre la tension $u(t)$ à ses bornes et le courant $i(t)$ qui le traverse.

On distingue notamment :
\begin{itemize}
  \item les dipôles passifs : résistance, condensateur, inductance ;
  \item les dipôles actifs : sources de tension et sources de courant ;
  \item les dipôles linéaires et non linéaires ;
  \item les dipôles récepteurs et générateurs selon leur échange de puissance avec le reste du circuit.
\end{itemize}

\subsection{Sources de tension et de courant}

Une source idéale de tension impose une tension indépendante du courant débité. Une source idéale de courant impose un courant indépendant de la tension à ses bornes.

\begin{center}
\begin{circuitikz}
  \draw (0,0) to[V,l=$E$] (0,3);
  \draw (3,0) to[I,l=$I_0$] (3,3);
  \node[below] at (0,0) {source de tension};
  \node[below] at (3,0) {source de courant};
\end{circuitikz}
\end{center}

Une source réelle de tension est modélisée par une source idéale en série avec une résistance interne. Une source réelle de courant est modélisée par une source idéale en parallèle avec une résistance interne.

\subsection{Énergie et puissance}

La puissance électrique instantanée est définie par :
\[
p(t)=u(t)i(t).
\]

Elle s'exprime en watts. Elle traduit la quantité d'énergie fournie ou consommée par unité de temps :
\[
W(t_1,t_2)=\int_{t_1}^{t_2}p(t)\,\mathrm{d}t.
\]

En régime continu, lorsque $U$ et $I$ sont constants :
\[
P=UI,
\qquad
W=P\Delta t.
\]

\subsection{Point de fonctionnement}

Chaque point de la caractéristique tension-courant d'un dipôle représente un fonctionnement possible. Le point de fonctionnement réel d'un circuit correspond à l'intersection de la caractéristique du générateur et de celle du récepteur. Il fixe simultanément la tension aux bornes du récepteur et le courant qui le traverse.

\subsection{Conventions récepteur et générateur}

Le courant et la tension sont des grandeurs algébriques. Leur signe dépend de l'orientation choisie pour les flèches.

\begin{itemize}
  \item En convention récepteur, les flèches de tension et de courant sont de sens opposés.
  \item En convention générateur, les flèches de tension et de courant sont de même sens.
\end{itemize}

La convention est choisie en fonction du fonctionnement attendu. Une résistance fonctionne toujours comme récepteur. Un condensateur ou une inductance peuvent alternativement recevoir puis restituer de l'énergie.

\subsection{La résistance}

La résistance obéit à la loi d'Ohm :
\[
u(t)=Ri(t),
\qquad
R\text{ en }\si{\ohm}.
\]

La conductance vaut $G=1/R$ et s'exprime en siemens :
\[
i(t)=Gu(t).
\]

\subsubsection{Associations de résistances}

Pour des résistances en série :
\[
R_{\mathrm{eq}}=R_1+R_2+\cdots+R_n.
\]

Pour des résistances en parallèle :
\[
\frac{1}{R_{\mathrm{eq}}}=\frac{1}{R_1}+\frac{1}{R_2}+\cdots+\frac{1}{R_n}.
\]

Dans le cas de deux résistances en parallèle :
\[
R_{\mathrm{eq}}=\frac{R_1R_2}{R_1+R_2}.
\]

La puissance dissipée par effet Joule est :
\[
P_R=Ri^2=\frac{u^2}{R}=ui.
\]

\subsection{Le condensateur}

Un condensateur stocke de l'énergie sous forme électrostatique. Sa relation courant-tension est :
\[
i(t)=C\frac{\mathrm{d}u(t)}{\mathrm{d}t}.
\]

La tension aux bornes d'un condensateur ne peut pas varier instantanément. L'énergie emmagasinée vaut :
\[
W_C=\frac{1}{2}Cu^2.
\]

En régime continu établi, le condensateur se comporte comme un circuit ouvert.

Pour des condensateurs en parallèle :
\[
C_{\mathrm{eq}}=C_1+C_2+\cdots+C_n.
\]

Pour des condensateurs en série :
\[
\frac{1}{C_{\mathrm{eq}}}=\frac{1}{C_1}+\frac{1}{C_2}+\cdots+\frac{1}{C_n}.
\]

\subsection{L'inductance}

Une inductance stocke de l'énergie sous forme magnétique. Sa relation courant-tension est :
\[
u(t)=L\frac{\mathrm{d}i(t)}{\mathrm{d}t}.
\]

Le courant dans une inductance ne peut pas varier instantanément. L'énergie emmagasinée vaut :
\[
W_L=\frac{1}{2}Li^2.
\]

En régime continu établi, une inductance idéale se comporte comme un court-circuit.

Pour des inductances non couplées en série :
\[
L_{\mathrm{eq}}=L_1+L_2+\cdots+L_n.
\]

Pour des inductances non couplées en parallèle :
\[
\frac{1}{L_{\mathrm{eq}}}=\frac{1}{L_1}+\frac{1}{L_2}+\cdots+\frac{1}{L_n}.
\]

\subsection{Quadrants de fonctionnement}

Les conventions de fléchage ne déterminent pas à elles seules le mode de fonctionnement du dipôle. Celui-ci dépend du signe de la puissance $p=ui$.

En convention récepteur :
\begin{itemize}
  \item si $p>0$, le dipôle reçoit de l'énergie : il fonctionne en récepteur ;
  \item si $p<0$, le dipôle fournit de l'énergie : il fonctionne en générateur.
\end{itemize}

En convention générateur :
\begin{itemize}
  \item si $p>0$, le dipôle fonctionne en générateur ;
  \item si $p<0$, il fonctionne en récepteur.
\end{itemize}

Pour une résistance, la caractéristique $u=Ri$ passe par l'origine. En convention récepteur, les quadrants 1 et 3 correspondent au fonctionnement récepteur. Pour une source de tension idéale représentée en convention générateur, les quadrants sont interprétés de manière inverse.

\section{Outils et lois de l'électrocinétique}\label{sec:elec-lois}

\subsection{Lois de Kirchhoff}

Un circuit peut être constitué d'un ensemble de dipôles associés. On définit :
\begin{description}
  \item[Nœud :] connexion entre au moins deux dipôles ;
  \item[Branche :] portion de circuit comprise entre deux nœuds consécutifs ;
  \item[Maille :] chemin fermé constitué de plusieurs branches.
\end{description}

\subsubsection{Loi des nœuds}

La somme algébrique des courants arrivant à un nœud est nulle :
\[
\sum_k i_k=0.
\]

On peut aussi écrire que la somme des courants entrants est égale à la somme des courants sortants.

\subsubsection{Loi des mailles}

Dans une maille orientée, la somme algébrique des tensions est nulle :
\[
\sum_k u_k=0.
\]

Le signe d'une tension dépend de l'orientation de sa flèche par rapport au sens choisi pour parcourir la maille.

\subsection{Diviseur de tension}

Le diviseur de tension s'applique à des résistances placées en série et parcourues par le même courant.

\TLPontDiviseur

Pour deux résistances :
\[
u_s=E\frac{R_2}{R_1+R_2}.
\]

Plus généralement, la tension aux bornes de $R_k$ dans une association série vaut :
\[
U_k=U\frac{R_k}{\sum_i R_i}.
\]

\subsection{Diviseur de courant}

Le diviseur de courant s'applique à des résistances en parallèle soumises à la même tension.

Pour deux résistances parcourues par les courants $I_1$ et $I_2$, avec $I=I_1+I_2$ :
\[
I_1=I\frac{R_2}{R_1+R_2},
\qquad
I_2=I\frac{R_1}{R_1+R_2}.
\]

En utilisant les conductances :
\[
I_k=I\frac{G_k}{\sum_i G_i}.
\]

\subsection{Théorème de superposition}

Dans un circuit linéaire comportant plusieurs sources indépendantes, la tension ou le courant dans une branche est la somme algébrique des contributions obtenues en ne conservant qu'une source active à la fois.

Pour neutraliser les autres sources :
\begin{itemize}
  \item une source idéale de tension est remplacée par un court-circuit ;
  \item une source idéale de courant est remplacée par un circuit ouvert.
\end{itemize}

Le théorème s'applique aux tensions et aux courants, mais pas directement aux puissances.

\subsection{Théorème de Millman}

Le théorème de Millman permet de déterminer le potentiel d'un nœud relié à plusieurs branches contenant chacune une source de tension et une résistance.

Pour un nœud de potentiel $V$ relié à des potentiels $E_k$ par des résistances $R_k$ :
\[
V=\frac{\displaystyle\sum_k \frac{E_k}{R_k}}
        {\displaystyle\sum_k \frac{1}{R_k}}.
\]

Avec les conductances $G_k=1/R_k$ :
\[
V=\frac{\sum_k G_kE_k}{\sum_k G_k}.
\]

Le signe de chaque terme dépend de l'orientation de la source dans la branche correspondante.

\subsection{Théorème de Kennelly}

Le théorème de Kennelly permet de transformer un réseau de trois résistances montées en triangle en un réseau équivalent en étoile, et réciproquement.

Pour le passage triangle $R_{AB}$, $R_{BC}$, $R_{CA}$ vers étoile $R_A$, $R_B$, $R_C$ :
\[
R_A=\frac{R_{AB}R_{CA}}{R_{AB}+R_{BC}+R_{CA}},
\]
\[
R_B=\frac{R_{AB}R_{BC}}{R_{AB}+R_{BC}+R_{CA}},
\]
\[
R_C=\frac{R_{BC}R_{CA}}{R_{AB}+R_{BC}+R_{CA}}.
\]

Pour le passage étoile vers triangle :
\[
R_{AB}=R_A+R_B+\frac{R_AR_B}{R_C},
\]
\[
R_{BC}=R_B+R_C+\frac{R_BR_C}{R_A},
\]
\[
R_{CA}=R_C+R_A+\frac{R_CR_A}{R_B}.
\]

\subsection{Synthèse électrocinétique}

Pour choisir rapidement une méthode de résolution :
\begin{itemize}
  \item une source avec une association série : diviseur de tension ;
  \item une source avec une association parallèle : diviseur de courant ;
  \item plusieurs sources de tension reliées à un même nœud : théorème de Millman ;
  \item un mélange de sources de tension et de courant dans un circuit linéaire : superposition ;
  \item un réseau résistif non réductible par série-parallèle : transformation de Kennelly ou lois de Kirchhoff.
\end{itemize}

\subsubsection{Petites applications}

À partir des circuits proposés dans le document de cours :
\begin{enumerate}
  \item déterminer l'expression d'un potentiel de nœud $V_1$ ;
  \item déterminer l'expression d'un potentiel de nœud $V_4$ ;
  \item en déduire la différence de potentiel aux bornes d'une résistance donnée ;
  \item justifier le choix de la méthode utilisée parmi Kirchhoff, diviseur, Millman et superposition.
\end{enumerate}

% -----------------------------------------------------------------------------
% Les parties suivantes sont encore provisoires. Elles seront remplacées dans
% l'ordre du document source lors des prochaines reprises.
% -----------------------------------------------------------------------------

\section{Amplificateur opérationnel}\label{sec:elec-aop}

Un amplificateur opérationnel idéal présente un gain différentiel infini, des courants d'entrée nuls et une résistance de sortie nulle. En régime linéaire avec contre-réaction négative,
\[
i_+=i_-=0
\qquad\text{et}\qquad
u_+=u_-.
\]

\subsection{Montage inverseur}
\TLAOPInverseur
\[
\boxed{\frac{u_s}{u_e}=-\frac{R_2}{R_1}}.
\]

\subsection{Montage non-inverseur}
\TLAOPNonInverseur
\[
\boxed{\frac{u_s}{u_e}=1+\frac{R_2}{R_1}}.
\]

\section{Filtrage analogique}\label{sec:elec-filtrage}

\subsection{Filtre RC passe-bas}
\TLFiltreRCPasseBas
\[
H(p)=\frac{1}{1+RCp},
\qquad
f_c=\frac{1}{2\pi RC}.
\]

\subsection{Filtre RC passe-haut}
\TLFiltreRCPasseHaut
\[
H(p)=\frac{RCp}{1+RCp}.
\]

\section{Échantillonnage}\label{sec:elec-echantillonnage}

\TLEchantillonnageSignal
\[
f_e=\frac{1}{T_e},
\qquad
\boxed{f_e>2f_{\max}}.
\]

\section{Conversion analogique-numérique}\label{sec:elec-can}

Pour un convertisseur sur $N$ bits et une plage $[U_{\min},U_{\max}]$,
\[
q=\frac{U_{\max}-U_{\min}}{2^N}.
\]

\section{Conversion numérique-analogique}\label{sec:elec-cna}

Un convertisseur numérique-analogique produit une tension ou un courant proportionnel au mot binaire appliqué en entrée.

\section{À retenir}\label{sec:elec-retenir}

\begin{itemize}
  \item Une chaîne d'acquisition associe capteur, conditionnement, filtrage, conversion et traitement.
  \item La nature analogique, logique ou numérique de l'information conditionne son traitement.
  \item Le choix d'un capteur repose sur son étendue, sa résolution, sa précision, sa sensibilité, sa rapidité et sa bande passante.
  \item Le conditionnement adapte le signal du capteur à l'unité de traitement.
  \item Les lois de Kirchhoff, les diviseurs, Millman, la superposition et Kennelly constituent les principaux outils de résolution des circuits résistifs.
\end{itemize}

